%=====================================================================
% Matriz de trazabilidad
%=====================================================================
\section*{Matriz de trazabilidad}
\addcontentsline{toc}{section}{Matriz de trazabilidad}

Documento de control. No se lee de corrido: se consulta para responder tres preguntas.

\begin{flujo}
  \item ¿Qué caso de uso justifica esta tabla? — \textbf{Matriz 1}
  \item ¿Qué se rompe si cambio este caso de uso? — \textbf{Matriz 2} y \textbf{Matriz 3}
  \item ¿Quedó alguna pantalla sin caso o alguna tabla sin quién la escriba? — \textbf{Matriz 1} y \textbf{Matriz 4}
\end{flujo}

Todos los casos de uso están redactados a detalle. La Matriz 1 se obtuvo directamente de la sección \textit{Datos que toca} de cada caso.

\subsection*{Matriz 1. Estructuras de datos y el caso de uso que las justifica}
\addcontentsline{toc}{subsection}{Matriz 1. Estructuras de datos y el caso de uso que las justifica}
Ninguna fila puede quedarse sin al menos un caso que la escriba. Una tabla que nadie escribe o no hace falta, o falta el caso de uso que la llena.

\begin{center}\small
\begin{longtable}{|L{0.240\textwidth}|L{0.280\textwidth}|L{0.240\textwidth}|L{0.100\textwidth}|}
\hline
\textbf{Estructura} & \textbf{Escrita por} & \textbf{Leída por} & \textbf{Estado} \\ \hline
\endfirsthead
\hline
\textbf{Estructura} & \textbf{Escrita por} & \textbf{Leída por} & \textbf{Estado} \\ \hline
\endhead
\textit{AceptacionTerminos} & \mbox{CU-02}, \mbox{CU-07} & — & correcta \\ \hline
\textit{Amistad} & \mbox{CU-11}, \mbox{CU-31}, \mbox{CU-33}, \mbox{CU-49} & \mbox{CU-16}, \mbox{CU-28}, \mbox{CU-30}, \mbox{CU-35} & correcta \\ \hline
\textit{Apelacion} & \mbox{CU-43}, \mbox{CU-45}, \mbox{CU-47} & — & correcta \\ \hline
\textit{Avatar} & \mbox{CU-11}, \mbox{CU-12} & \mbox{CU-15}, \mbox{CU-16} & correcta \\ \hline
\textit{BitacoraAcceso} & \mbox{CU-01}, \mbox{CU-02}, \mbox{CU-03}, \mbox{CU-04}, \mbox{CU-05}, \mbox{CU-06}, \mbox{CU-07}, \mbox{CU-09}, \mbox{CU-10}, \mbox{CU-11} & \mbox{CU-48} & correcta \\ \hline
\textit{BitacoraModeracion} & \mbox{CU-43}, \mbox{CU-44}, \mbox{CU-46}, \mbox{CU-47}, \mbox{CU-48} & — & correcta \\ \hline
\textit{Bloqueo} & \mbox{CU-11}, \mbox{CU-33} & \mbox{CU-16}, \mbox{CU-17}, \mbox{CU-19}, \mbox{CU-28}, \mbox{CU-30}, \mbox{CU-31}, \mbox{CU-32}, \mbox{CU-34}, \mbox{CU-35} & correcta \\ \hline
\textit{Caja} & \mbox{CU-25}, \mbox{CU-39} & \mbox{CU-40} & correcta \\ \hline
\textit{CajaContenido} & \mbox{CU-39} & \mbox{CU-40} & correcta \\ \hline
\textit{CodigoSegundoFactor} & \mbox{CU-01} & — & correcta \\ \hline
\textit{ColaEmparejamiento} & \mbox{CU-11}, \mbox{CU-17}, \mbox{CU-33} & \mbox{CU-18}, \mbox{CU-19}, \mbox{CU-27}, \mbox{CU-32} & correcta \\ \hline
\textit{Division} & — \textit{(catálogo, \mbox{D-09})} & \mbox{CU-15}, \mbox{CU-16}, \mbox{CU-35} & precargada \\ \hline
\textit{EnlaceEspectador} & \mbox{CU-25}, \mbox{CU-34} & — & correcta \\ \hline
\textit{EnlaceRed} & \mbox{CU-11}, \mbox{CU-12} & \mbox{CU-15}, \mbox{CU-16} & correcta \\ \hline
\textit{Equipamiento} & \mbox{CU-02}, \mbox{CU-06}, \mbox{CU-08}, \mbox{CU-12}, \mbox{CU-14} & \mbox{CU-01}, \mbox{CU-15}, \mbox{CU-16}, \mbox{CU-17}, \mbox{CU-18} & correcta \\ \hline
\textit{EstadisticaModo} & \mbox{CU-02}, \mbox{CU-06}, \mbox{CU-24}, \mbox{CU-25} & \mbox{CU-01}, \mbox{CU-15}, \mbox{CU-16}, \mbox{CU-17}, \mbox{CU-18}, \mbox{CU-31}, \mbox{CU-35} & correcta \\ \hline
\textit{HistorialDivision} & \mbox{CU-25} & \mbox{CU-15} & correcta \\ \hline
\textit{HistorialNickname} & \mbox{CU-13} & — & correcta \\ \hline
\textit{Invitacion} & \mbox{CU-11}, \mbox{CU-32}, \mbox{CU-33} & — & correcta \\ \hline
\textit{Jugada} & \mbox{CU-20}, \mbox{CU-21}, \mbox{CU-27} & \mbox{CU-25}, \mbox{CU-26}, \mbox{CU-34}, \mbox{CU-37}, \mbox{CU-43} & correcta \\ \hline
\textit{Leccion} & — \textit{(catálogo, \mbox{D-09})} & \mbox{CU-41} & precargada \\ \hline
\textit{Mensaje} & \mbox{CU-18}, \mbox{CU-28} & \mbox{CU-19}, \mbox{CU-43} & correcta \\ \hline
\textit{Modo} & \mbox{CU-16} & \mbox{CU-02}, \mbox{CU-06}, \mbox{CU-17}, \mbox{CU-18}, \mbox{CU-19}, \mbox{CU-21}, \mbox{CU-22}, \mbox{CU-23}, \mbox{CU-24}, \mbox{CU-25}, \mbox{CU-27}, \mbox{CU-36}, \mbox{CU-37}, \mbox{CU-40} & correcta \\ \hline
\textit{MotivoReporte} & — \textit{(catálogo, \mbox{D-09})} & \mbox{CU-42}, \mbox{CU-43} & precargada \\ \hline
\textit{MovimientoMoneda} & \mbox{CU-25}, \mbox{CU-38}, \mbox{CU-39} & \mbox{CU-40} & correcta \\ \hline
\textit{NivelIA} & — \textit{(catálogo, \mbox{D-09})} & \mbox{CU-27}, \mbox{CU-36} & precargada \\ \hline
\textit{ObjetoCosmetico} & — \textit{(catálogo, \mbox{D-09})} & \mbox{CU-02}, \mbox{CU-06}, \mbox{CU-08}, \mbox{CU-14}, \mbox{CU-38}, \mbox{CU-39}, \mbox{CU-40} & precargada \\ \hline
\textit{OfertaTablas} & \mbox{CU-22}, \mbox{CU-23}, \mbox{CU-24} & \mbox{CU-26} & correcta \\ \hline
\textit{PalabraProhibida} & \mbox{CU-45} & \mbox{CU-02}, \mbox{CU-07}, \mbox{CU-13}, \mbox{CU-28} & correcta \\ \hline
\textit{Participacion} & \mbox{CU-16}, \mbox{CU-17}, \mbox{CU-18}, \mbox{CU-20}, \mbox{CU-21}, \mbox{CU-23}, \mbox{CU-24}, \mbox{CU-25}, \mbox{CU-26}, \mbox{CU-27}, \mbox{CU-44} & \mbox{CU-01}, \mbox{CU-05}, \mbox{CU-10}, \mbox{CU-11}, \mbox{CU-15}, \mbox{CU-19}, \mbox{CU-22}, \mbox{CU-28}, \mbox{CU-29}, \mbox{CU-30}, \mbox{CU-32}, \mbox{CU-34}, \mbox{CU-36}, \mbox{CU-37} & correcta \\ \hline
\textit{ParticipacionObjeto} & \mbox{CU-18}, \mbox{CU-27} & \mbox{CU-26}, \mbox{CU-29}, \mbox{CU-34}, \mbox{CU-37} & correcta \\ \hline
\textit{Partida} & \mbox{CU-16}, \mbox{CU-17}, \mbox{CU-18}, \mbox{CU-20}, \mbox{CU-21}, \mbox{CU-25}, \mbox{CU-27} & \mbox{CU-15}, \mbox{CU-22}, \mbox{CU-23}, \mbox{CU-24}, \mbox{CU-26}, \mbox{CU-34}, \mbox{CU-36}, \mbox{CU-37}, \mbox{CU-40}, \mbox{CU-42}, \mbox{CU-43} & correcta \\ \hline
\textit{ProgresoTutorial} & \mbox{CU-41} & — & correcta \\ \hline
\textit{Ranura} & — \textit{(catálogo, \mbox{D-09})} & \mbox{CU-02}, \mbox{CU-06}, \mbox{CU-14} & precargada \\ \hline
\textit{Reporte} & \mbox{CU-11}, \mbox{CU-42}, \mbox{CU-43}, \mbox{CU-44}, \mbox{CU-47} & \mbox{CU-37}, \mbox{CU-46} & correcta \\ \hline
\textit{Sala} & \mbox{CU-18}, \mbox{CU-25}, \mbox{CU-32} & \mbox{CU-19}, \mbox{CU-34} & correcta \\ \hline
\textit{SalaParticipante} & \mbox{CU-11}, \mbox{CU-18}, \mbox{CU-19}, \mbox{CU-25}, \mbox{CU-32}, \mbox{CU-33} & \mbox{CU-27} & correcta \\ \hline
\textit{Sancion} & \mbox{CU-43}, \mbox{CU-44} & \mbox{CU-01}, \mbox{CU-04}, \mbox{CU-17}, \mbox{CU-18}, \mbox{CU-19}, \mbox{CU-28}, \mbox{CU-45}, \mbox{CU-46} & correcta \\ \hline
\textit{Sesion} & \mbox{CU-01}, \mbox{CU-03}, \mbox{CU-05}, \mbox{CU-06}, \mbox{CU-07}, \mbox{CU-09}, \mbox{CU-10}, \mbox{CU-11}, \mbox{CU-44} & \mbox{CU-08}, \mbox{CU-12}, \mbox{CU-13}, \mbox{CU-14}, \mbox{CU-15}, \mbox{CU-16}, \mbox{CU-17}, \mbox{CU-18}, \mbox{CU-19}, \mbox{CU-20}, \mbox{CU-21}, \mbox{CU-22}, \mbox{CU-23}, \mbox{CU-24}, \mbox{CU-26}, \mbox{CU-27}, \mbox{CU-28}, \mbox{CU-30}, \mbox{CU-31}, \mbox{CU-32}, \mbox{CU-33}, \mbox{CU-34}, \mbox{CU-35}, \mbox{CU-36}, \mbox{CU-37}, \mbox{CU-38}, \mbox{CU-39}, \mbox{CU-40}, \mbox{CU-41}, \mbox{CU-42}, \mbox{CU-43}, \mbox{CU-45}, \mbox{CU-46}, \mbox{CU-47}, \mbox{CU-48}, \mbox{CU-49} & correcta \\ \hline
\textit{Silencio} & \mbox{CU-11}, \mbox{CU-33} & \mbox{CU-16}, \mbox{CU-28} & correcta \\ \hline
\textit{Solicitud} & \mbox{CU-11}, \mbox{CU-30}, \mbox{CU-31}, \mbox{CU-33} & \mbox{CU-16} & correcta \\ \hline
\textit{TipoCaja} & — \textit{(catálogo, \mbox{D-09})} & \mbox{CU-39}, \mbox{CU-40} & precargada \\ \hline
\textit{TipoCajaObjeto} & — \textit{(catálogo, \mbox{D-09})} & \mbox{CU-39} & precargada \\ \hline
\textit{TokenRecuperacion} & \mbox{CU-03}, \mbox{CU-09} & — & correcta \\ \hline
\textit{TokenVerificacionCorreo} & \mbox{CU-01}, \mbox{CU-02}, \mbox{CU-04}, \mbox{CU-07}, \mbox{CU-09} & — & correcta \\ \hline
\textit{Usuario} & \mbox{CU-01}, \mbox{CU-02}, \mbox{CU-03}, \mbox{CU-04}, \mbox{CU-06}, \mbox{CU-07}, \mbox{CU-08}, \mbox{CU-09}, \mbox{CU-11}, \mbox{CU-12}, \mbox{CU-13}, \mbox{CU-25}, \mbox{CU-38}, \mbox{CU-39}, \mbox{CU-43}, \mbox{CU-44}, \mbox{CU-46}, \mbox{CU-47} & \mbox{CU-14}, \mbox{CU-15}, \mbox{CU-16}, \mbox{CU-30}, \mbox{CU-31}, \mbox{CU-32}, \mbox{CU-33}, \mbox{CU-34}, \mbox{CU-35}, \mbox{CU-36}, \mbox{CU-37}, \mbox{CU-40}, \mbox{CU-42}, \mbox{CU-48} & correcta \\ \hline
\textit{UsuarioObjeto} & \mbox{CU-02}, \mbox{CU-06}, \mbox{CU-25}, \mbox{CU-38}, \mbox{CU-39} & \mbox{CU-08}, \mbox{CU-12}, \mbox{CU-14} & correcta \\ \hline
\end{longtable}
\end{center}

\textbf{Sin huérfanas.} Las 47 estructuras que aparecen en los casos tienen al menos un caso de uso que las escribe, salvo los 10 catálogos, que por la decisión \mbox{D-09} se cargan con \texttt{insertar\_\allowbreak{}datos\_\allowbreak{}prueba.sql} y se mantienen por SQL. Esa es una elección, no un olvido, y conviene que aparezca escrita en el documento de entrega.

\textbf{Estructuras sin tabla propia.} Tres conceptos que podrían parecer tablas se resuelven sin ellas: el historial de elo sale de las \textit{Participacion} clasificatorias (\mbox{CU-15}), el marcador entre dos jugadores se calcula al consultarlo (\mbox{CU-16}) y los espectadores se llevan en memoria (\mbox{CU-34}). Cada caso justifica la decisión en su observación.

\subsection*{Matriz 2. Dependencias entre casos de uso}
\addcontentsline{toc}{subsection}{Matriz 2. Dependencias entre casos de uso}
Un caso \textit{depende} de otro cuando lo invoca, lo continúa o presupone su resultado. Cambiar el de la derecha obliga a revisar el de la izquierda.

\begin{center}\small
\begin{longtable}{|L{0.155\textwidth}|L{0.170\textwidth}|L{0.555\textwidth}|}
\hline
\textbf{Caso} & \textbf{Depende de} & \textbf{Por qué} \\ \hline
\endfirsthead
\hline
\textbf{Caso} & \textbf{Depende de} & \textbf{Por qué} \\ \hline
\endhead
\mbox{CU-01} & \mbox{CU-04} & Sin verificación, ninguna cuenta llega a \texttt{ACTIVA} y el \mbox{FA-09} rechaza a todas \\ \hline
\mbox{CU-01} & \mbox{CU-44} & Lee las sanciones de ámbito de cuenta \\ \hline
\mbox{CU-03} & \mbox{CU-01} & Nace de su \mbox{FA-03} \\ \hline
\mbox{CU-04} & \mbox{CU-02}, \mbox{CU-07} & Consume el token que cualquiera de los dos generó \\ \hline
\mbox{CU-05} & \mbox{CU-26} & Su \mbox{FA-01} ofrece volver a la partida \\ \hline
\mbox{CU-07} & \mbox{CU-06} & Solo puede vincular lo que el \mbox{CU-06} creó \\ \hline
\mbox{CU-07} & \mbox{CU-04} & Deja la cuenta en \texttt{PENDIENTE}, que solo el \mbox{CU-04} resuelve \\ \hline
\mbox{CU-12} & \mbox{CU-14} & El título es un equipamiento: comparten operación \\ \hline
\mbox{CU-13} & \mbox{CU-42} & \textit{HistorialNickname} existe para que un reportado no se vuelva irreconocible \\ \hline
\mbox{CU-14} & \mbox{CU-38}, \mbox{CU-39} & Solo se equipa lo que se posee \\ \hline
\mbox{CU-17} & \mbox{CU-25} & Lee el elo que solo el \mbox{CU-25} escribe \\ \hline
\mbox{CU-17} & \mbox{CU-26} & Su \mbox{FA-01} ofrece volver a la partida en curso \\ \hline
\mbox{CU-20}, \mbox{CU-21} & \mbox{CU-17} & Sin partida creada no hay turno \\ \hline
\mbox{CU-20}, \mbox{CU-21} & \mbox{CU-25} & Le entregan el control al detectar el término \\ \hline
\mbox{CU-25} & \mbox{CU-20}, \mbox{CU-21} & Cuenta los muros leyendo sus \textit{Jugada} \\ \hline
\mbox{CU-28} & \mbox{CU-44} & Lee las sanciones de ámbito de chat \\ \hline
\mbox{CU-30} & \mbox{CU-31} & La \textit{Amistad} la crea el otro al aceptar \\ \hline
\mbox{CU-38} & \mbox{CU-25} & Las monedas solo las otorga el fin de partida \\ \hline
\mbox{CU-39} & \mbox{CU-38} & Comparten \textit{UsuarioObjeto} y \textit{MovimientoMoneda} \\ \hline
\mbox{CU-42} & \mbox{CU-28}, \mbox{CU-20}, \mbox{CU-21} & Adjunta por referencia lo que ellos escribieron \\ \hline
\mbox{CU-44} & \mbox{CU-42}, \mbox{CU-43} & Resuelve el reporte que ellos crean y revisan \\ \hline
\mbox{CU-44} & \mbox{CU-25} & Cierra la participación del sancionado y le cede el control \\ \hline
\mbox{CU-44} & \mbox{CU-46} & Exige el rol que la administración otorga \\ \hline
\mbox{CU-08} & \mbox{CU-01}, \mbox{CU-06} & Se ofrece al terminar el primer acceso de cualquiera de los dos \\ \hline
\mbox{CU-09} & \mbox{CU-04} & El cambio de correo se confirma con el enlace del \mbox{CU-04} (\mbox{FA-07}) \\ \hline
\mbox{CU-11} & \mbox{CU-01} & La recuperación dentro de los catorce días ocurre al iniciar sesión (\mbox{FA-17}) \\ \hline
\mbox{CU-19}, \mbox{CU-32} & \mbox{CU-18} & Entrar con código o por invitación supone una sala creada, y el inicio de la partida lo hace el \mbox{CU-18} \\ \hline
\mbox{CU-22}, \mbox{CU-23}, \mbox{CU-24} & \mbox{CU-25} & Son formas de término; el cierre y la puntuación los hace el \mbox{CU-25} \\ \hline
\mbox{CU-24} & \mbox{CU-26} & Reclamar la victoria convierte en abandono una desconexión que el \mbox{CU-26} no resolvió \\ \hline
\mbox{CU-27} & \mbox{CU-25} & La experiencia y el repaso de errores se calculan en su \mbox{FA-01} \\ \hline
\mbox{CU-34} & \mbox{CU-37} & Al terminar la partida o caducar el enlace se ofrece la repetición \\ \hline
\mbox{CU-43} & \mbox{CU-44}, \mbox{CU-45} & Resuelve los reportes que acaban en sanción y las apelaciones \\ \hline
\mbox{CU-44}, \mbox{CU-43} & \mbox{CU-46} & Exigen el rol que otorga la administración \\ \hline
\mbox{CU-49} & \mbox{CU-31} & Busca la amistad por el mismo par ordenado con que se guardó \\ \hline
\end{longtable}
\end{center}

\subsection*{Matriz 3. Decisiones y los casos que las aplican}
\addcontentsline{toc}{subsection}{Matriz 3. Decisiones y los casos que las aplican}
Cambiar una decisión obliga a reescribir flujos, no solo tablas. Esta es la lista de a quién hay que avisar.

\begin{center}\small
\begin{longtable}{|L{0.308\textwidth}|L{0.592\textwidth}|}
\hline
\textbf{Decisión} & \textbf{Casos afectados} \\ \hline
\endfirsthead
\hline
\textbf{Decisión} & \textbf{Casos afectados} \\ \hline
\endhead
\mbox{D-01} · Varias sesiones, una partida & \mbox{CU-01}, \mbox{CU-05}, \mbox{CU-17}, \mbox{CU-25}, \mbox{CU-10}, \mbox{CU-18}, \mbox{CU-19}, \mbox{CU-27}, \mbox{CU-32} \\ \hline
\mbox{D-02} · Invitado como \textit{Usuario} & \mbox{CU-01}, \mbox{CU-02}, \mbox{CU-06}, \mbox{CU-07}, \mbox{CU-13}, \mbox{CU-30}, \mbox{CU-38}, \mbox{CU-42} \\ \hline
\mbox{D-03} · Seis ranuras & \mbox{CU-02}, \mbox{CU-06}, \mbox{CU-12}, \mbox{CU-14}, \mbox{CU-17}, \mbox{CU-38}, \mbox{CU-39} \\ \hline
\mbox{D-04} · Elo por modo & \mbox{CU-02}, \mbox{CU-06}, \mbox{CU-17}, \mbox{CU-25}, \mbox{CU-15}, \mbox{CU-16}, \mbox{CU-35} \\ \hline
\mbox{D-05} · Sanciones con ámbito & \mbox{CU-01}, \mbox{CU-04}, \mbox{CU-17}, \mbox{CU-28}, \mbox{CU-44}, \mbox{CU-45} \\ \hline
\mbox{D-06} · Registro mínimo & \mbox{CU-02}, \mbox{CU-07} \\ \hline
\mbox{D-07} · Borrado lógico & \mbox{CU-03}, \mbox{CU-04}, \mbox{CU-30}, \mbox{CU-42}, \mbox{CU-11} \\ \hline
\mbox{D-08} · Sin revisión de avatares & \mbox{CU-12} \\ \hline
\mbox{D-09} · Catálogos precargados & \mbox{CU-02}, \mbox{CU-07}, \mbox{CU-13}, \mbox{CU-14}, \mbox{CU-28}, \mbox{CU-38}, \mbox{CU-39}, \mbox{CU-42} \\ \hline
\mbox{D-10} · Emotes sin persistir & \mbox{CU-29} \\ \hline
\mbox{D-11} · Una participación, incluida la IA & \mbox{CU-17}, \mbox{CU-18}, \mbox{CU-19}, \mbox{CU-27} \\ \hline
\mbox{D-12} · Partidas contra la IA en el servidor & \mbox{CU-25}, \mbox{CU-27}, \mbox{CU-36}, \mbox{CU-37} \\ \hline
\mbox{D-13} · Primera vez: peón e icono & \mbox{CU-01}, \mbox{CU-02}, \mbox{CU-06}, \mbox{CU-08} \\ \hline
\mbox{D-14} · Tablas solo entre dos & \mbox{CU-22}, \mbox{CU-25} \\ \hline
\mbox{D-15} · Abandono y desconexión & \mbox{CU-24}, \mbox{CU-25}, \mbox{CU-26} \\ \hline
\mbox{D-16} · Salas privadas & \mbox{CU-18}, \mbox{CU-19}, \mbox{CU-32} \\ \hline
\mbox{D-17} · Recuperación y anonimización & \mbox{CU-01}, \mbox{CU-11} \\ \hline
\mbox{D-18} · Tutorial en la cuenta & \mbox{CU-41} \\ \hline
\mbox{D-19} · Espectadores sin persistencia & \mbox{CU-25}, \mbox{CU-34} \\ \hline
\mbox{D-20} · Cambio de correo por token & \mbox{CU-04}, \mbox{CU-09} \\ \hline
\end{longtable}
\end{center}

\subsection*{Matriz 4. Prototipos y el caso de uso que los cubre}
\addcontentsline{toc}{subsection}{Matriz 4. Prototipos y el caso de uso que los cubre}
Solo los prototipos con caso asignado. Los descartados están en la sección de decisiones de diseño, con su motivo.

\begin{center}\small
\begin{longtable}{|L{0.443\textwidth}|L{0.457\textwidth}|}
\hline
\textbf{Prototipos} & \textbf{Caso} \\ \hline
\endfirsthead
\hline
\textbf{Prototipos} & \textbf{Caso} \\ \hline
\endhead
6c & \mbox{CU-01}, \mbox{CU-06} \\ \hline
6d & \mbox{CU-02}, \mbox{CU-07} \\ \hline
9a, 9b & \mbox{CU-03} \\ \hline
9e & \mbox{CU-06}, \mbox{CU-07} \\ \hline
9f & \mbox{CU-05}, \mbox{CU-09}, \mbox{CU-10} \\ \hline
9g & \mbox{CU-11} \\ \hline
9c, 18d & \mbox{CU-12}, \mbox{CU-13} \\ \hline
9d & \mbox{CU-13} \\ \hline
4a, 4b, 4c, 17e & \mbox{CU-14} \\ \hline
5a, 16c, 16d & \mbox{CU-15} \\ \hline
5b & \mbox{CU-16} \\ \hline
16a, 16b, 1f, 1g, 7d & \mbox{CU-17} \\ \hline
1h, 5d, 5e & \mbox{CU-18}, \mbox{CU-19} \\ \hline
17c, 8c, 8d, 20a, 7a, 17d & \mbox{CU-20} \\ \hline
15a, 15b, 21a, 20b, 7b, 19b & \mbox{CU-21} \\ \hline
19d & \mbox{CU-22}, \mbox{CU-23} \\ \hline
8e & \mbox{CU-24}, \mbox{CU-25} \\ \hline
2g, 2h, 11a, 11b, 15c, 15d, 12g & \mbox{CU-25} \\ \hline
6e, 6f & \mbox{CU-26} \\ \hline
8f, 12f & \mbox{CU-27} \\ \hline
20c, 5f, 1b, 10g & \mbox{CU-28} \\ \hline
12e & \mbox{CU-29} \\ \hline
2a, 2b, 5c, 13b, 13c & \mbox{CU-30}, \mbox{CU-31}, \mbox{CU-32} \\ \hline
10a, 10c & \mbox{CU-33} \\ \hline
10f, 10g, 13a, 13d, 13e & \mbox{CU-34} \\ \hline
2c, 2d & \mbox{CU-35} \\ \hline
2e, 17f & \mbox{CU-36} \\ \hline
2f, 15e & \mbox{CU-37} \\ \hline
4d, 4e, 4f & \mbox{CU-38} \\ \hline
14a, 14b, 14c, 14d, 14e, 18c & \mbox{CU-39} \\ \hline
16e & \mbox{CU-40} \\ \hline
6a, 6b & \mbox{CU-41} \\ \hline
10b & \mbox{CU-42} \\ \hline
10d, 10e & \mbox{CU-44}, \mbox{CU-45} \\ \hline
\end{longtable}
\end{center}

\subsection*{Matriz 5. Reglas transversales}
\addcontentsline{toc}{subsection}{Matriz 5. Reglas transversales}
Reglas que aparecen con el mismo contenido en varios casos. Si cambia una, hay que cambiarla en todos los sitios de su fila, porque están numeradas de forma local a cada caso  y nada las mantiene sincronizadas por sí solo.

\begin{center}\small
\begin{longtable}{|L{0.445\textwidth}|L{0.455\textwidth}|}
\hline
\textbf{Regla} & \textbf{Dónde aparece} \\ \hline
\endfirsthead
\hline
\textbf{Regla} & \textbf{Dónde aparece} \\ \hline
\endhead
Las contraseñas y los códigos nunca viajan ni se guardan en claro & \mbox{CU-01} \mbox{RN-02}, \mbox{CU-02} \mbox{RN-03}, \mbox{CU-03} \mbox{RN-02}, \mbox{CU-04} \mbox{RN-06}, \mbox{CU-07} \mbox{RN-02} \\ \hline
Toda validación del cliente se repite en el servidor & \mbox{CU-02} \mbox{RN-10}, \mbox{CU-07} \mbox{RN-09}, \mbox{CU-12} \mbox{RN-07}, \mbox{CU-13} \mbox{RN-08}, \mbox{CU-20} \mbox{RN-07}, \mbox{CU-21} \mbox{RN-08}, \mbox{CU-28} \mbox{RN-08} \\ \hline
El servidor es la única autoridad sobre la partida & \mbox{CU-20} \mbox{RN-01}, \mbox{CU-21} \mbox{RN-01}, \mbox{CU-25} \mbox{RN-04} \\ \hline
El cliente pide el estado en lugar de reenviar & \mbox{CU-20} \mbox{RN-13}, \mbox{CU-21} \mbox{RN-13}, \mbox{CU-28} \mbox{RN-10}, \mbox{CU-38} \mbox{RN-09} \\ \hline
El \texttt{nickname} es único y se compara sin acentos ni mayúsculas & \mbox{CU-02} \mbox{RN-01}, \mbox{CU-07}, \mbox{CU-13} \mbox{RN-02} \\ \hline
Todo movimiento de monedas deja fila y el saldo es su suma & \mbox{CU-25} \mbox{RN-08}, \mbox{CU-38} \mbox{RN-05}, \mbox{CU-39} \\ \hline
Un \textit{Usuario} tiene una \textit{Ranura} equipada de cada tipo, siempre & \mbox{CU-02}, \mbox{CU-06} \mbox{RN-06}, \mbox{CU-14} \mbox{RN-01} \\ \hline
Una sola \textit{Participacion} sin terminar por jugador & \mbox{CU-05} \mbox{RN-04}, \mbox{CU-17} \mbox{RN-01}, \mbox{CU-25} \mbox{RN-09} \\ \hline
El filtro de palabras se aplica en el servidor y no dice qué término saltó & \mbox{CU-02} \mbox{RN-09}, \mbox{CU-07}, \mbox{CU-13} , \mbox{CU-28} \mbox{RN-02} \\ \hline
Todo lo cosmético no altera las reglas del juego & \mbox{CU-14} \mbox{RN-05}, \mbox{CU-38} \mbox{RN-06} \\ \hline
\end{longtable}
\end{center}

\subsection*{Lo que esta matriz deja al descubierto}
\addcontentsline{toc}{subsection}{Lo que esta matriz deja al descubierto}
\begin{flujo}
  \item \textbf{El \mbox{CU-25} es el cuello de botella del modelo.} Escribe en once estructuras y es el único origen del elo, de las monedas, de la división y de los contadores. Si algo se prueba con datos reales, que sea este.
  \item \textbf{\textit{Usuario} y \textit{Sesion} tienen muchos escritores.} Conviene que las restricciones de sus columnas —dominios de \texttt{estado\_\allowbreak{}cuenta}, \texttt{rol} y \texttt{motivo\_\allowbreak{}cierre}— se declaren en la base de datos con \texttt{CHECK}, y no dependan de que cada caso respete los valores.
  \item \textbf{\texttt{saldo\_\allowbreak{}monedas} es una redundancia controlada.} Lo escriben tres casos y lo verifica un cuarto (\mbox{CU-40}). Debe justificarse en la evidencia de tercera forma normal.
  \item \textbf{Las reglas transversales no tienen dueño.} La numeración local por caso es cómoda para leer un caso aislado, pero deja reglas repetidas en varios sitios sin nada que las mantenga iguales. La Matriz 5 es el único lugar donde se ven juntas.
\end{flujo}
